\cleardoublepage
\phantomsection              % for hyperlinks/bookmarks
\chapter*{\Huge\bfseries\hfill\textRL{اهداء}}
\thispagestyle{noheader}
\addcontentsline{toc}{chapter}{\textRL{اهداء}}

\vspace{0.5cm}
\begin{RLtext}

بسم اللّه الرحمان الرحيم،  
و الصلاة و السلام على أشرف المرسلين.  

الحمد للّه الذي ما تمّ جهدٌ ولا خُتم سعيٌ إلا بفضله،  
فبجزيل عطائه ارتقيت لهذه المرحلة،  
و بتوفيق منه سبحانه تكلّل جهدي و اجتهادي لعيش هذه اللحظة؛  
لحظة التخرج من رحلتي التعليمية الحافلة.  

إلى أهل غزة المستضعفين،  
في حين يتمحور موضوعي حول عنصر أساسي في التغذية،  
يعانون هم اليوم من المجاعة والإبادة، فاللهم انصرهم وفك كربهم.  

لمن حملتني وهنًا على وهن،  
ولم تتوانَ يومًا عن دعمي والتضحية لأجلي...  
إلى أمي الغالية التي طالما رأت سعادتها في سعادتي.  

لمن علّمني الصبر والتفاؤل والإحسان،  
وشجعني في كل خطوة خطوتها،  
واحتفى بكل إنجاز صغير لي وكأنه شيء عظيم...  
لأبي الذي لم يرجُ مني سوى نجاحي.  

إلى روح عمّي الطيبة الطاهرة،  
لازالت ذكراك في قلبي حين رجوت نجاحي وتخرجي.  
رحمة الله عليك.  

لإخوتي، أحبتي، وسندي:  
إيمان، هشام، سندس، وعلي...  
وجودكم بجانبي قوة أعتز بها.  
أسأل الله أن يديمكم بقربي دائمًا، ويفتح عليكم بكل ما هو جميل.  

لزميلتي التي شاركتني الكفاح في هذا المشروع،  
من سهرت معي الليالي، وشاركتني المخاوف والضحكات،  
إليك يا زوليخة صديقتي الغالية، ونعم الصحبة الصالحة.  

إلى رفيقاتي اللواتي وقفن بجانبي دائمًا:  
خديجة رفيقة الروح، غادة طيبة القلب،  
سمية الصديق الصدوق، وسهى الرقيقة اللطيفة.  

لعائلتي الثانية، عائلتي الترجيحية،  
التي ضمّتني لحضنها وكونتني لنفع هذه الأمة.  

وأخيرًا... إلى نفسي التي آمنت بي،  
بإذن الله سأغدو يومًا كما أردت، بتوفيق من الله ربي.  

\hfill أميرة

\end{RLtext}

\pagebreak
